\documentclass[12pt]{article}

\usepackage{amsmath, amssymb}
\usepackage{geometry}
\usepackage{parskip}
\usepackage{booktabs}
\usepackage{array}
\usepackage{listings}
\usepackage{forest}
\usepackage{verbatim}

\geometry{margin=1in}

\title{OurPL Mini Assignment}
\author{
Greyson Knippel \\[4pt]
CPSC326 --- Gonzaga University \\[4pt]
Professor Johnson \\[4pt]
\textit{Compiled with \LaTeX}
}
\date{\today}

\begin{document}

\maketitle

% ─────────────────────────────────────────────
\section*{answers!!!}
% ─────────────────────────────────────────────


\subsection*{1. is it left recursive?}

no. it doesn't derive itself on the very first symbol.

it avoids it by using the (...)* pattern

\subsection*{2. is it left factorable?}

no, because it has no common shared prefixes.

the \texttt{(...)*} loop structure already eliminates conflicting prefixes

\subsection*{3. is it ambiguous?}

no. there aren't multuiple ways to derive a string here, you can only derive it one way through a primary.

% ─────────────────────────────────────────────
\section*{3 example derivations}
% ─────────────────────────────────────────────

\textbf{number}, \textbf{bool}, and \textbf{string}.

\subsection*{derivation 1 --- number: \texttt{2 + 2 / 2}}

\begin{align*}
&\langle\textit{expression}\rangle \\
\Rightarrow\;&\langle\textit{equality}\rangle \\
\Rightarrow\;&\langle\textit{comparison}\rangle \\
\Rightarrow\;&\langle\textit{term}\rangle \\
\Rightarrow\;&\langle\textit{factor}\rangle\;(+\;\langle\textit{factor}\rangle)^* \\
\Rightarrow\;&\langle\textit{unary}\rangle\;(+\;\langle\textit{factor}\rangle)^* \\
\Rightarrow\;&\langle\textit{primary}\rangle\;(+\;\langle\textit{factor}\rangle)^* \\
\Rightarrow\;&2\;(+\;\langle\textit{factor}\rangle)^* \\
\Rightarrow\;&2\;+\;\langle\textit{factor}\rangle \\
\Rightarrow\;&2\;+\;\langle\textit{unary}\rangle\;(/\;\langle\textit{unary}\rangle)^* \\
\Rightarrow\;&2\;+\;\langle\textit{primary}\rangle\;(/\;\langle\textit{unary}\rangle)^* \\
\Rightarrow\;&2\;+\;2\;(/\;\langle\textit{unary}\rangle)^* \\
\Rightarrow\;&2\;+\;2\;/\;\langle\textit{unary}\rangle \\
\Rightarrow\;&2\;+\;2\;/\;\langle\textit{primary}\rangle \\
\Rightarrow\;&2\;+\;2\;/\;2
\end{align*}


\subsection*{derivation 2 --- bool: \texttt{"dog" == "zebra"}}

\begin{align*}
&\langle\textit{expression}\rangle \\
\Rightarrow\;&\langle\textit{equality}\rangle \\
\Rightarrow\;&\langle\textit{comparison}\rangle\;(==\;\langle\textit{comparison}\rangle)^* \\
\Rightarrow\;&\langle\textit{term}\rangle\;(==\;\langle\textit{comparison}\rangle)^* \\
\Rightarrow\;&\langle\textit{factor}\rangle\;(==\;\langle\textit{comparison}\rangle)^* \\
\Rightarrow\;&\langle\textit{unary}\rangle\;(==\;\langle\textit{comparison}\rangle)^* \\
\Rightarrow\;&\langle\textit{primary}\rangle\;(==\;\langle\textit{comparison}\rangle)^* \\
\Rightarrow\;&\texttt{"dog"}\;(==\;\langle\textit{comparison}\rangle)^* \\
\Rightarrow\;&\texttt{"dog"}\;==\;\langle\textit{comparison}\rangle \\
\Rightarrow\;&\texttt{"dog"}\;==\;\langle\textit{term}\rangle \\
\Rightarrow\;&\texttt{"dog"}\;==\;\langle\textit{factor}\rangle \\
\Rightarrow\;&\texttt{"dog"}\;==\;\langle\textit{unary}\rangle \\
\Rightarrow\;&\texttt{"dog"}\;==\;\langle\textit{primary}\rangle \\
\Rightarrow\;&\texttt{"dog"}\;==\;\texttt{"zebra"}
\end{align*}

\subsection*{derivation 3 --- string: \texttt{"flower"}}

\begin{align*}
&\langle\textit{expression}\rangle \\
\Rightarrow\;&\langle\textit{equality}\rangle \\
\Rightarrow\;&\langle\textit{comparison}\rangle \\
\Rightarrow\;&\langle\textit{term}\rangle \\
\Rightarrow\;&\langle\textit{factor}\rangle \\
\Rightarrow\;&\langle\textit{unary}\rangle \\
\Rightarrow\;&\langle\textit{primary}\rangle \\
\Rightarrow\;&\texttt{"flower"}
\end{align*}

% ─────────────────────────────────────────────
\section*{parsing trees}
% ─────────────────────────────────────────────


\subsection*{parse tree 1: \texttt{"Hello" + " World" == "Hello World"}}

\begin{align*}
&\langle\textit{expression}\rangle \\
\Rightarrow\;&\langle\textit{equality}\rangle \\
\Rightarrow\;&\langle\textit{comparison}\rangle\;(==\;\langle\textit{comparison}\rangle)^* \\
\Rightarrow\;&\langle\textit{term}\rangle\;(==\;\langle\textit{comparison}\rangle)^* \\
\Rightarrow\;&\langle\textit{factor}\rangle\;(+\;\langle\textit{factor}\rangle)^*\;(==\;\langle\textit{comparison}\rangle)^* \\
\Rightarrow\;&\langle\textit{unary}\rangle\;(+\;\langle\textit{factor}\rangle)^*\;(==\;\langle\textit{comparison}\rangle)^* \\
\Rightarrow\;&\langle\textit{primary}\rangle\;(+\;\langle\textit{factor}\rangle)^*\;(==\;\langle\textit{comparison}\rangle)^* \\
\Rightarrow\;&\texttt{"Hello"}\;(+\;\langle\textit{factor}\rangle)^*\;(==\;\langle\textit{comparison}\rangle)^* \\
\Rightarrow\;&\texttt{"Hello"}\;+\;\langle\textit{factor}\rangle\;(==\;\langle\textit{comparison}\rangle)^* \\
\Rightarrow\;&\texttt{"Hello"}\;+\;\langle\textit{unary}\rangle\;(==\;\langle\textit{comparison}\rangle)^* \\
\Rightarrow\;&\texttt{"Hello"}\;+\;\langle\textit{primary}\rangle\;(==\;\langle\textit{comparison}\rangle)^* \\
\Rightarrow\;&\texttt{"Hello"}\;+\;\texttt{" World"}\;(==\;\langle\textit{comparison}\rangle)^* \\
\Rightarrow\;&\texttt{"Hello"}\;+\;\texttt{" World"}\;==\;\langle\textit{comparison}\rangle \\
\Rightarrow\;&\texttt{"Hello"}\;+\;\texttt{" World"}\;==\;\langle\textit{term}\rangle \\
\Rightarrow\;&\texttt{"Hello"}\;+\;\texttt{" World"}\;==\;\langle\textit{factor}\rangle \\
\Rightarrow\;&\texttt{"Hello"}\;+\;\texttt{" World"}\;==\;\langle\textit{unary}\rangle \\
\Rightarrow\;&\texttt{"Hello"}\;+\;\texttt{" World"}\;==\;\langle\textit{primary}\rangle \\
\Rightarrow\;&\texttt{"Hello"}\;+\;\texttt{" World"}\;==\;\texttt{"Hello World"}
\end{align*}

\textbf{LLk value:} 1. you only need to look ahead one token to see whats next.

\subsection*{parse tree 2: \texttt{1 + 2 * 3 - 4}}

\begin{align*}
&\langle\textit{expression}\rangle \\
\Rightarrow\;&\langle\textit{equality}\rangle \\
\Rightarrow\;&\langle\textit{comparison}\rangle \\
\Rightarrow\;&\langle\textit{term}\rangle \\
\Rightarrow\;&\langle\textit{factor}\rangle\;(+\;|\;-\;\langle\textit{factor}\rangle)^* \\
\Rightarrow\;&\langle\textit{unary}\rangle\;(+\;|\;-\;\langle\textit{factor}\rangle)^* \\
\Rightarrow\;&\langle\textit{primary}\rangle\;(+\;|\;-\;\langle\textit{factor}\rangle)^* \\
\Rightarrow\;&1\;(+\;|\;-\;\langle\textit{factor}\rangle)^* \\
\Rightarrow\;&1\;+\;\langle\textit{factor}\rangle\;(-\;\langle\textit{factor}\rangle)^* \\
\Rightarrow\;&1\;+\;\langle\textit{unary}\rangle\;(*\;\langle\textit{unary}\rangle)^*\;(-\;\langle\textit{factor}\rangle)^* \\
\Rightarrow\;&1\;+\;\langle\textit{primary}\rangle\;(*\;\langle\textit{unary}\rangle)^*\;(-\;\langle\textit{factor}\rangle)^* \\
\Rightarrow\;&1\;+\;2\;(*\;\langle\textit{unary}\rangle)^*\;(-\;\langle\textit{factor}\rangle)^* \\
\Rightarrow\;&1\;+\;2\;*\;\langle\textit{unary}\rangle\;(-\;\langle\textit{factor}\rangle)^* \\
\Rightarrow\;&1\;+\;2\;*\;\langle\textit{primary}\rangle\;(-\;\langle\textit{factor}\rangle)^* \\
\Rightarrow\;&1\;+\;2\;*\;3\;(-\;\langle\textit{factor}\rangle)^* \\
\Rightarrow\;&1\;+\;2\;*\;3\;-\;\langle\textit{factor}\rangle \\
\Rightarrow\;&1\;+\;2\;*\;3\;-\;\langle\textit{unary}\rangle \\
\Rightarrow\;&1\;+\;2\;*\;3\;-\;\langle\textit{primary}\rangle \\
\Rightarrow\;&1\;+\;2\;*\;3\;-\;4
\end{align*}

\textbf{LLk value:} 1 token again.

\subsection*{parse tree 3: \texttt{"Storms, these steamed cremlings are divine!" + " Rock, you've done it again."}}

\begin{align*}
&\langle\textit{expression}\rangle \\
\Rightarrow\;&\langle\textit{equality}\rangle \\
\Rightarrow\;&\langle\textit{comparison}\rangle \\
\Rightarrow\;&\langle\textit{term}\rangle \\
\Rightarrow\;&\langle\textit{factor}\rangle\;(+\;\langle\textit{factor}\rangle)^* \\
\Rightarrow\;&\langle\textit{unary}\rangle\;(+\;\langle\textit{factor}\rangle)^* \\
\Rightarrow\;&\langle\textit{primary}\rangle\;(+\;\langle\textit{factor}\rangle)^* \\
\Rightarrow\;&\texttt{"Storms\ldots divine!"}\;(+\;\langle\textit{factor}\rangle)^* \\
\Rightarrow\;&\texttt{"Storms\ldots divine!"}\;+\;\langle\textit{factor}\rangle \\
\Rightarrow\;&\texttt{"Storms\ldots divine!"}\;+\;\langle\textit{unary}\rangle \\
\Rightarrow\;&\texttt{"Storms\ldots divine!"}\;+\;\langle\textit{primary}\rangle \\
\Rightarrow\;&\texttt{"Storms\ldots divine!"}\;+\;\texttt{" Rock\ldots again."}
\end{align*}

\textbf{LLk value:} 1 yet again :D

% ─────────────────────────────────────────────
\section*{is this grammar LL(1)?}
% ─────────────────────────────────────────────

yep. as shown above, all trees only have to look ahead once to know where to go.
it doesnt recurse left, and doesnt have left factoring which is bad.


% ─────────────────────────────────────────────
\section*{operator precedence}
% ─────────────────────────────────────────────

\subsection*{how does this grammar enforce order of operations?}

by whichever operator is deepest, it determines the order of the operation.

for example, / and * are evaluated last because they are the deepest, compared to + and - which
aren't the deepest. and you have to evaluate the leaves before the rest of the tree.

\subsection*{can we add a string to a number?}

yes. there's nothing preventing us from doing so. this also just isn't the parser's
job. it has no derivation rules to stop us from simply adding a string to a number when deriving term.

\end{document}